\section{Sub-Agents}
\label{sec:subagents}

Each agent $a\!\in\!\mathcal{A}$ is a self-describing module rather
than a hand-coded pipeline stage: the orchestrator discovers,
schedules, and executes agents through a uniform descriptor interface,
with no agent-specific switches in the orchestration code.

\paragraph{Descriptor.}
Every sub-agent ships a single \texttt{SubAgentDescriptor} carrying
its planner-facing \texttt{Inputs}/\texttt{Output}/\texttt{Purpose} block,
a \texttt{build\_input}\,$:$\,$(\textit{step\_id},\,\textit{resolved\_artifacts})\!\to\!\textit{TypedInput}$
function, factory hooks for the LLM body and the evaluator, and an
optional materializer factory for media-producing agents (image,
video, audio, composition). Adding an agent is a single sub-package
plus one line in \texttt{AGENT\_REGISTRY}.

\paragraph{Execution lifecycle.}
For each \texttt{PlanStep} $=(a,\textit{step\_id})$ Assistant runs:
InputResolver matches workspace artifacts to the agent's
\texttt{[label]}s, \texttt{build\_input} packs them into a Pydantic
typed input, the agent's LLM emits structured JSON, the evaluator gates
$L_1$ (rule-based) and $L_2$ (LLM creative), the materializer (if any)
generates binaries, and $L_3$ checks the produced assets. Failure
emits \texttt{rework\_notes} that retry \texttt{generate()}; failure
beyond budget surfaces \texttt{FAILED} and the Director picks among
retry / replan-tail / skip. A separate rejection hatch lets an agent
abort with a structured upstream-rejection instead of fabricating
output.

\paragraph{Quality gate.}
The evaluator owns three output-only layers (no cross-validation
against upstream): $L_1$ deterministic structural rules
(ID format and sequencing, e.g.\ \texttt{shot\_id}\,$\!=\!$\,
\texttt{sh\_NNN} globally sequential, \texttt{keyframe\_count}\,$\!=\!1$
per shot, metric consistency); $L_2$ LLM-as-judge over agent-specific
creative dimensions; $L_3$ post-materialization asset checks (URI
success rate, format, assembly). The same evaluator generates the
\texttt{rework\_notes} fed back to \texttt{generate()}, so the gate
and the self-correction signal share one definition.

\paragraph{Strict-out, loose-in.}
A sub-agent enforces its own output schema strictly (Pydantic $+$ $L_1$
$+$ evaluator) but reads upstream artifacts with no key-by-name
assumptions: \texttt{build\_input} dumps each upstream payload as raw
JSON text into a \texttt{*\_json\_text:\,str} field, and the agent's
LLM reads the text and infers the upstream structure from prompt
context. Cross-agent invariants (ID formats, per-shot counts,
sequentiality) are the producer's obligation --- enforced in its
prompt and its $L_1$ check, surfaced as rework on drift --- never
patched silently in the consumer. This keeps the inter-agent contract
a one-line method signature instead of an
\textit{O}$(N\!\times\!M)$ web of hard-coded field accesses across $20$
descriptors.

\paragraph{Inventory.}
$|\mathcal{A}|\!=\!20$ across modalities (text, image, video, audio,
mix, composition); the per-modality breakdown is given in
\S\ref{sec:bench}. There is no fixed DAG over $\mathcal{A}$: any
prefix / re-ordering is a valid plan as long as each agent's input
labels are satisfiable at execution time. Routing is therefore a
learnable problem (\S\ref{sec:sft}, \S\ref{sec:grpo}) over a closed
but combinatorially-large plan space, and adding a new agent does not
require retraining the router beyond extending its catalog.
